% !TeX spellcheck = ru_RU
% !TEX root = vkr.tex

\section*{Введение}
\thispagestyle{withCompileDate}

% Формат из 4х частей рекомендуется в курсе Д.~Кознова~\cite{koznov} по написанию текстов.

% \begin{enumerate}
%     \item Известная информация (background/обзор).
%     \item Неизвестная информация (пробел в знаниях, \enquote{Gap}).
%     \item Гипотезы, вопросы, цели~--- \enquote{что болит}, что будет решать Ваша работа.
%     \item Подход, план решения задачи, предлагаемое решение.
% \end{enumerate}

% Последний абзац должен читаться и быть понятен в отрыве от других трёх.
% Никакие абзацы нумеровать нельзя.

% Части (абзацы) должны занять максимум две страницы, идеально уложиться в одну.

% С.-П. Джонс~\cite{SPJGreatPaper} предлагает несколько другой формат написания введения.
% Вполне возможно, что если Ваша работа про языки программирования, то его формат будет удачнее.

% Введение и заключение читают чаще всего, поэтому они должны быть \enquote{вылизаны} до блеска.

% \blfootnote{
%     Иногда рецензенту полезно знать какого числа компилировался текст, чтобы оценить актуальность версии текста.
%     В этом случае полезно вставлять в текст дату сборки.
%     Для совсем официальных релизов документа это не вполне канон.\\
%     Также здесь имеет смысл указать, если работа сделана на деньги, например, Российского Фонда Фундаментальных Исследований (РФФИ) по гранту номер такой-то, и т.п.}

Строительная область включает в себя большое количество процессов, каждый из которых требует контроля и анализа результатов.
Эффективность работы напрямую зависит от качества контроля. Основной подход к контролю над объектами --- это использование специально-обученных людей,
таких как охранники, прорабы, инженеры и другие специалисты. В данном варианте есть несколько проблем: во-первых, любая ручная работа требует денежных затрат, 
что для долгого объекта может навредить бюджету, во-вторых, невозможность повлиять на человеческий фактор:
сокрытие информации, некачественная работа, воровство и т.д., а в третьих, результаты контроля хотят получать разные люди, 
от руководителя до заказчиков объекта.

В настоящее время многие компании стараются автоматизировать процессы и переложить большинство задач программным системам.
Развитие технологий, методов машинного обучения и компьютерного зрения позволяет автоматизировать контролем над строительными объектами.
Существуют много решений, в ключающие в себя турникеты, умные системы видеонаблюдения, датчики движения и другие подходы.
Зачастую такие системы являются дорогостоящими и не все компании могут себе позволить их внедрение.

Возникает потребность в открытых умных системах для контроля над объектами, внедрение которых будет дешёвым, 
а результаты контроля будут доступны для заинтересованных лиц. 
Большим плюсом для таких систем является возможность их интегрирования в существующее оборудование, 
возможность модернизации и расширение функционала.

В данной работе будет разработана открытая система сбора статистики со строительных объектов.
В качестве источника данных будет использоваться RTSP стрим с камеры видеонаблюдения, охватывающая зону выхода / входа,
после будет проивзодиться обработка данных и генерация отчёта со статистикой с дальнешей рассылкой в удобный месседжер.